% ============================================================
%  Section 9: Epistemic Uncertainty Analysis
%  H&E to Sirius Red Virtual Staining Paper
% ============================================================

\section{Epistemic Uncertainty Analysis}
\label{sec:uncertainty}

The scaling study~(Section~\ref{sec:results_scaling}) identifies how model capacity and
data volume jointly affect translation quality, but it treats each configuration as a
deterministic predictor.
In practice, an unpaired I2I model trained from a single random initialisation offers no
indication of where its predictions are reliable and where they should be treated with
caution.
We address this gap by training deep ensembles for the best-performing configuration
from each model family---where \textit{best} is judged solely by task-based CPA
agreement, for reasons detailed in Section~\ref{sec:ens_selection}---and constructing
spatially resolved epistemic uncertainty maps that quantify the tile-level confidence
of each model.

% -------------------------------------------------------
\subsection{Model Selection for Ensemble Analysis}
\label{sec:ens_selection}
% -------------------------------------------------------

For each of the six model families (CycleGAN, UNIT, MUNIT, DCLGAN, UVCGAN, MIUDiff),
we select the configuration from the 54-run scaling study that achieves the best
CPA agreement on the held-out test set (Section~\ref{sec:task_eval}), yielding six
configurations that enter the ensemble analysis.

\paragraph{Rationale for CPA-only selection.}
A central finding of this paper is that standard perceptual and distributional
metrics---FID (Section~\ref{sec:fid}), LPIPS (Section~\ref{sec:lpips}), and
patch-based SSIM (Section~\ref{sec:ssim})---are unreliable proxies for the clinical
utility of virtual staining in an unpaired setting
(Section~\ref{sec:results_scaling}).
Selecting configurations by a composite of these metrics would implicitly endorse
their validity for model selection, contradicting this conclusion.
We therefore select exclusively on CPA agreement, the only metric that directly
measures whether a virtual SR image supports the same downstream collagen
quantification as a real SR image.
This choice is methodologically consistent: the configurations that enter the
uncertainty analysis are those judged best by the criterion we argue is the
appropriate one, not by the metrics we argue are insufficient.

\paragraph{Selection procedure.}
All 54 configurations are ranked by CPA agreement from 1 (best) to 54 (worst).
Within each model family, the configuration with the best (lowest-ranked) CPA
agreement score is selected as the representative for ensemble training.
In the event of a tie within a family, the configuration with the lower FID is
preferred as a secondary tiebreak.

\paragraph{Selected configurations.}
\textcolor{red}{[TODO: insert table of selected configurations once 54-run results are
available. For each model family, report the winning size (Small / Medium / Large),
data fraction (25\% / 50\% / 100\%), CPA agreement score, and CPA rank out of 54.
Include FID, LPIPS, and patch-SSIM for these configurations as supplementary
columns to illustrate the relationship---or divergence---between task-based and
perceptual rankings.]}

\paragraph{Ensemble training.}
For each selected configuration, ten models are trained independently from distinct
random seeds, using the identical hyperparameters and data as the corresponding
scaling-study run.
Using the same fixed data ensures that variance in ensemble member outputs reflects
genuine epistemic uncertainty arising from model initialisation rather than exposure to
different training examples.
The ensemble size of $K{=}10$ is motivated by convergence analysis showing that
calibration errors plateau at approximately ten members for deep ensembles in
biomedical image segmentation~\cite{siddiqui2024uncertainty}, and is consistent
with prior ensemble-based uncertainty work in virtual staining
applications~\cite{trustworthy_cell_ensemble2023}.

% -------------------------------------------------------
\subsection{Uncertainty Map Construction}
\label{sec:unc_maps}
% -------------------------------------------------------

Given an ensemble of $K{=}10$ independently trained generators
$\{G^{(k)}_{A \to B}\}_{k=1}^{K}$ for a selected configuration, we derive a per-tile
mean prediction and a spatially resolved epistemic uncertainty map for each H\&E test
tile $x \in \mathcal{D}_{\text{test}}$.

\paragraph{Ensemble mean prediction.}
The ensemble mean image $\bar{y}(x)$ is computed as the pixel-wise arithmetic mean of
the ten generator outputs:
\begin{equation}
  \bar{y}(x) = \frac{1}{K} \sum_{k=1}^{K} G^{(k)}_{A \to B}(x).
\end{equation}
This mean image serves as the primary virtual stain prediction for each
ensemble-evaluated configuration and is used in all downstream metric computations
described in Section~\ref{sec:metrics}.

\paragraph{Epistemic uncertainty map.}
Epistemic uncertainty at each pixel location $i$ is estimated as the pixel-wise
variance across ensemble members:
\begin{equation}
  \sigma^2_c(x, y) = \frac{1}{K-1} \sum_{k=1}^{K}
    \bigl(G^{(k)}_{A \to B}(x)_{c,x,y} - \bar{y}(x)_{c,x,y}\bigr)^2,
    \quad c \in \{R, G, B\}.
\end{equation}
The three per-channel variance maps are summed to obtain a single-channel scalar
uncertainty map:
\begin{equation}
  U(x, y) = \sigma^2_R(x,y) + \sigma^2_G(x,y) + \sigma^2_B(x,y),
\end{equation}
giving $\mathbf{U}(x) \in \mathbb{R}^{H \times W}$, where $H{\times}W = 256{\times}256$
pixels.
Using the sample variance~($K{-}1$ denominator) is consistent with the unbiased
estimator of the population variance under the finite ensemble~\cite{lakshminarayanan2017simple}.
Summing rather than averaging channels preserves the total disagreement signal: a
pixel where all three channels disagree should register higher uncertainty than one
where only one channel varies.
High values in $\mathbf{U}(x)$ indicate pixels where the ensemble members disagree
substantially, reflecting regions where the model is uncertain about the correct
translation---for example, due to ambiguous texture, rare tissue morphologies, or
artefacts not well represented in the training distribution.

\paragraph{Visualisation.}
For qualitative inspection, each test tile is displayed as a three-panel composite: the
input H\&E tile, the ensemble mean virtual SR prediction $\bar{y}(x)$, and the
uncertainty map $\mathbf{U}(x)$ rendered as a heatmap (magma colorscale, normalised
globally to $[0,1]$ using dataset-wide $p_1$--$p_{99}$ percentile bounds for
comparability across tiles) overlaid on the mean prediction at 50\% opacity.
Tile-level scalar uncertainty is summarised as $\bar{\sigma}^2(x) =
\frac{1}{HW}\sum_i \sigma^2_i(x)$, the spatial mean of $\mathbf{U}(x)$, and used in
the correlation analyses of Section~\ref{sec:unc_vs_error}.
Qualitative uncertainty maps are shown for representative tiles spanning low,
intermediate, and high $\bar{\sigma}^2(x)$ values for each of the six model families,
enabling direct visual comparison of uncertainty structure across architectures.


% -------------------------------------------------------
\subsection{Uncertainty Calibration}
\label{sec:unc_calibration}
% -------------------------------------------------------

A spatially resolved uncertainty map is only useful if it is \textit{calibrated}:
when the model assigns high uncertainty to a pixel, the translation should be
measurably worse at that pixel.
Verifying this requires an independent per-pixel error signal.
In a supervised setting this would be the absolute deviation from a known target;
in the unpaired setting of this work no such target exists, because the H\&E and SR
images originate from different physical sections.
We therefore use \textit{cycle-reconstruction error} as a proxy for per-pixel
translation difficulty.
MC~dropout~\cite{gal2016dropout} was not adopted as an alternative to deep ensembles
because its uncertainty estimates require extensive dropout-rate tuning to achieve
calibration, and miscalibration persists even after tuning in complex
tasks~\cite{siddiqui2024uncertainty,trustworthy_i2i2025}; deep ensembles yield
well-calibrated estimates with minimal hyperparameter
sensitivity~\cite{lakshminarayanan2017simple,siddiqui2024uncertainty}.

\paragraph{Error proxy: judge-based cycle-reconstruction MAE.}
For a translated tile $B' = G_{A \to B}(A)$, the cycle-reconstruction error at pixel
$(x, y)$ is defined as
\begin{equation}
  E(x, y) = \bigl|A(x, y) - F_{\mathrm{judge}}(B')(x, y)\bigr|,
\end{equation}
where $F_{\mathrm{judge}}\colon B \to A$ is a fixed external GAN inverter applied
identically to every architecture's translated output.
We adopt this \textit{judge-based} variant in preference to the \textit{self-cycle}
variant, in which each model inverts its own translation.
Self-cycle error can be artificially low when the forward and inverse generators share
the same systematic bias: if both directions ignore the same tissue feature, the round
trip reconstructs the source faithfully despite a poor forward translation.
The judge-based variant removes this confound because the judge was trained
independently and cannot be jointly biased with the model under evaluation.
It additionally provides a calibration number for MIUDiff, which is a
unidirectional model ($A \to B$ only) with no learned inverse.
A single \textcolor{red}{[TODO: judge model architecture and checkpoint]} checkpoint is
frozen and reused as the judge for all six model families.
Calibration results are therefore stated as ``calibration with respect to
judge-based cycle-reconstruction error'' rather than against true translation error,
which is unavailable in this unpaired setting.

\paragraph{Tile preparation and tissue masking.}
Calibration is computed on a version of the test set re-tiled with zero overlap
(stride equal to tile size, $256{\times}256$ pixels), avoiding the double-counting
of border pixels that would bias across-tile and dataset-level statistics under
overlapping tiling.
All per-pixel metrics are restricted to tissue pixels identified by the preprocessing
tissue mask; background pixels --- where both $U$ and $E$ are near-zero --- are
excluded to prevent spurious inflation of rank correlations.

\paragraph{Calibration metrics.}
Once each tile yields a paired map $(U(x,y),\, E(x,y))$ over tissue pixels, we
evaluate calibration along four complementary dimensions following established
practice in predictive uncertainty
evaluation~\cite{kendall2017uncertainties,ilg2018ause,poggi2020sparsification}.

\subparagraph{Within-tile Spearman $\rho$ (spatial calibration).}
For each tile, $U$ and $E$ are flattened over tissue pixels and their Spearman rank
correlation $\rho$ is computed.
A value near $+1$ indicates that high-uncertainty pixels co-locate with high-error
pixels; near $0$ means the spatial structure of uncertainty is uninformative; a
negative value indicates anti-calibration.
We report mean~$\pm$~std of $\rho$ across all test tiles as the headline spatial
calibration figure.
Spearman rank correlation is preferred over Pearson because $U$ and $E$ are on
different scales and only their ordinal relationship is of interest.

\subparagraph{Across-tile Spearman and Pearson $\rho$ (tile-level calibration).}
Each tile is reduced to two scalars --- $\overline{U}$ and $\overline{E}$, the spatial
means over tissue pixels --- and both Spearman and Pearson correlations are computed
across all tiles.
This catches the failure mode where a model is spatially well-calibrated within each
tile but assigns identical mean uncertainty everywhere, making it useless for
flagging unreliable tiles for pathologist review.
Both coefficients are reported because tile-mean values tend to be more linearly
related than per-pixel values.

\subparagraph{AUSE --- sparsification error.}
The \textit{Area Under the Sparsification Error}
curve~\cite{ilg2018ause,poggi2020sparsification} measures whether uncertainty
correctly ranks pixels by their error magnitude, weighting rank inversions by how
much error they cost.
Pixels are sorted by predicted uncertainty in descending order; the mean error on the
retained fraction is tracked as pixels are progressively removed, giving the
\textit{predicted sparsification curve} $S_{\mathrm{pred}}(k)$.
An oracle curve $S_{\mathrm{oracle}}(k)$ is obtained by sorting on actual error $E$.
AUSE is the area between the two curves:
\begin{equation}
  \mathrm{AUSE} = \int_0^1 \bigl[S_{\mathrm{pred}}(k) - S_{\mathrm{oracle}}(k)
  \bigr]\,\mathrm{d}k.
\end{equation}
AUSE $= 0$ indicates that predicted and oracle orderings coincide; AUSE equal to the
random-ordering baseline indicates that uncertainty carries no information about error.
AUSE is computed per tile and averaged across tiles.
Unlike Spearman $\rho$, AUSE weights rank inversions by error magnitude, so a model
that correctly identifies the highest-error pixels scores well even if it
misranks low-error pixels.

\subparagraph{ECE and reliability diagram (calibration shape).}
Tissue pixels are binned by $U$-quantile into $B{=}10$ equal-mass bins.
Both the uncertainty and error maps are jointly rescaled to $[0,1]$ using
dataset-global $p_1$--$p_{99}$ bounds before binning, making the diagonal of the
reliability diagram a meaningful reference.
Within each bin $b$, the mean normalised uncertainty $\overline{U_n}_b$ and mean
normalised error $\overline{E_n}_b$ are computed and plotted against each other
(reliability diagram); perfect calibration corresponds to the identity diagonal.
An ECE-like summary score quantifies the weighted deviation from the
diagonal~\cite{guo2017ece}:
\begin{equation}
  \mathrm{ECE} = \sum_{b=1}^{B} \frac{n_b}{N}
    \bigl|\overline{U_n}_b - \overline{E_n}_b\bigr|,
\end{equation}
where $n_b$ is the number of pixels in bin $b$ and $N$ is the total.
This construction adapts the classification ECE of~\cite{guo2017ece} to the
regression setting by normalising both axes to $[0,1]$; it should not be compared
directly to ECE values from classification contexts.
Dataset-global normalisation bounds make ECE comparable across architectures within
this dataset.
Lower ECE indicates better agreement between uncertainty magnitude and error magnitude,
complementing the rank-based AUSE.

\paragraph{Caveats.}
The four metrics characterise the same underlying property from complementary angles
and should be read as a converging picture rather than independent evidence.
Cycle-reconstruction error is a proxy, not a ground truth: a perfect generator pair
with an imperfect judge will exhibit non-zero cycle error even for correct
translations, and cross-section biological differences between the H\&E and SR
sections add an irreducible noise floor to AUSE and ECE that no model can overcome.
These limitations are discussed further in Section~\ref{sec:discussion}.

% ============================================================
%  New bibliography entries for this section
%  (merge with references.bib)
% ============================================================
%
% @article{trustworthy_cell_ensemble2023,
%   author  = {Rauber, Paul and others},
%   title   = {Trustworthy in silico cell labeling via ensemble-based image translation},
%   journal = {Nature Communications},
%   year    = {2023},
%   note    = {PMC10663640},
%   url     = {https://pmc.ncbi.nlm.nih.gov/articles/PMC10663640/}
% }
%
% (Shared entries already listed:
%  lakshminarayanan2017simple, gal2016dropout, trustworthy_i2i2025)
